\documentclass[11pt,a4paper]{article}
\usepackage[T1]{fontenc}
\usepackage{lmodern}
\usepackage[utf8]{inputenc}
\usepackage[margin=2.5cm]{geometry}
\usepackage{booktabs}
\usepackage{array}
\usepackage{enumitem}
\usepackage{microtype}
\usepackage{titlesec}
\titlelabel{\thetitle.\quad}

\begin{document}

\setcounter{section}{1}
\section{Related Work}

This chapter describes the approaches from the literature to which the proposed study is compared. They are grouped by what they contribute to the keep-versus-switch decision: threshold rules used in practice, MDP models of routing, and the theory of switching costs. The chapter ends with a comparison table, the differences and similarities with the proposed approach, and the metrics where similar or better results are expected.

\subsection{Threshold Rules for Rerouting in Practice}

The classical way to avoid rerouting is to keep every flow on one path: ECMP hashes a flow to one of several equal-cost paths~\cite{hopps2000}, and routing weights are optimized offline for the expected traffic~\cite{fortz2000} or for a set of uncertain demands~\cite{applegate2003}. These methods do not react to changes. Reactive systems add a rule that moves traffic when it becomes worthwhile. Hedera detects large flows and places them again~\cite{alfares2010}, and CONGA balances load with feedback about congestion~\cite{alizadeh2014}.

\textbf{FlowBender}~\cite{kabbani2014} is the clearest example of a threshold rule. A flow changes its path when the fraction of ECN-marked packets exceeds a threshold $T$ for $N$ consecutive round-trip times. The cost of a switch is packet reordering, and the paper discusses the stability of the rule: flows can oscillate between paths, and a new path can overlap with the old one. The rule is evaluated by simulation and on a testbed, but it is not compared with an optimal policy.

\textbf{Kikuchi et al.}~\cite{kikuchi2025} decide in a software-defined network when a new packet flow allocation is worth installing. The reconfiguration overhead is the number of links with route changes, and the allocation is changed only when the ratio between the improvement in congestion and the overhead exceeds a threshold $\rho$. The rule is evaluated on the GEANT and US backbone networks against benchmark policies. The overhead depends on the route structure, which is close to the overlap penalty used here, but the rule is not compared with an optimal policy either.

The idea of acting only when a condition is violated also exists in control theory, where event-triggered control updates the controller only when the state leaves an allowed region~\cite{heemels2012}. It gives the name ``event-triggered'' to the problem studied here.

\subsection{MDP Models of Routing Decisions}

\textbf{Anjali et al.}~\cite{anjali2002} model the setup and resizing of label switched paths in MPLS networks as a continuous-time MDP with a discounted cost, and compute the optimal policy by value iteration. They also study a sub-optimal policy with a threshold (least one-step cost) and evaluate it by simulation. This is the closest model to the one used here, with the same comparison between an optimal and a threshold policy. The costs are given by the setup of paths and their resizing and they are not based on how much two paths share.

\textbf{Hamouda}~\cite{hamouda2020} models a node with two paths as an MDP with a switching cost, proves that the optimal policy has a state-dependent threshold structure, and compares it with a heuristic that uses a few static thresholds. On 100 random cases the heuristic loses 1.40\% on average and 4.29\% at most. This gives a reference value for the gap of a threshold policy. The switching cost is constant.

\textbf{Levering et al.}~\cite{levering2022} model routing under stochastically varying conditions as a stochastic shortest path problem, with a Markovian background process for congestion. They solve it by dynamic programming (value iteration) and propose two faster Dijkstra-like algorithms, whose distance to the optimum and run time are compared with the dynamic programming solution and with deterministic policies on road network data. The route can change at every intersection, and the cost of a change is not modeled. Their congestion model is the closest to the Markov-modulated costs of Chapter~3.

The classical references for these models are Puterman~\cite{puterman1994} and Bertsekas~\cite{bertsekas}, on which the value iteration used in this work is based.

\subsection{Switching Costs in Sequential Decision Making}

The theory of bandits with switching costs studies the same trade-off: staying on the current option, or paying a cost to change. Banks and Sundaram show that the Gittins index is no longer optimal when switching costs exist~\cite{banks1994}, and Asawa and Teneketzis propose an index for switching penalties~\cite{asawa1996}. Ni\~no-Mora gives a faster index algorithm based on the restless-bandit reformulation of Whittle~\cite{whittle1988} and measures the gap between the index policy and the optimal policy, which is below about 0.3\% in his experiments~\cite{ninomora2008}. These works explain why a simple rule is not optimal in general and give a way of measuring how much it loses. Their costs are for starting and stopping an option, and not for the structure of a route.

\subsection{Comparison}

\begin{table}[ht]
\centering
\small
\renewcommand{\arraystretch}{1.2}
\begin{tabular}{@{}>{\raggedright\arraybackslash}p{2.2cm}>{\raggedright\arraybackslash}p{2.4cm}>{\raggedright\arraybackslash}p{2.2cm}>{\raggedright\arraybackslash}p{2.6cm}>{\raggedright\arraybackslash}p{2.3cm}>{\raggedright\arraybackslash}p{2.2cm}@{}}
\toprule
Work & Decision & Switching cost & Policy & Compared with optimum & Evaluation \\
\midrule
Anjali et al.~\cite{anjali2002} & LSP setup and resizing & Setup cost & Optimal (MDP) and threshold & Yes & Simulation \\
Kikuchi et al.~\cite{kikuchi2025} & Flow reconfiguration & Links with route changes & Threshold on gain / overhead & No (benchmark policies) & GEANT, US backbone \\
Kabbani et al.~\cite{kabbani2014} & Flow path change & Packet reordering & Threshold with $N$ consecutive RTTs & No & Simulation, testbed \\
Hamouda~\cite{hamouda2020} & Admission and routing, 2 paths & Constant & Optimal and static thresholds & Yes (1.40\% mean, 4.29\% max) & 100 random cases \\
Levering et al.~\cite{levering2022} & Route at each intersection & Not modeled & Optimal (DP), fast algorithms & Yes & Road network data \\
Ni\~no-Mora~\cite{ninomora2008} & Bandits with switching costs & Start and stop costs & Index policy & Yes ($<$ 0.3\%) & Computational study \\
\midrule
This work & Keep or switch among $k$ paths & Constant and overlap-based & Threshold, hysteresis, overlap-aware & Yes, exact & Synthetic, then GEANT \\
\bottomrule
\end{tabular}
\caption{Comparison of the studied works with the proposed approach.}
\label{tab:comparison}
\end{table}

\subsection{Differences and Similarities with the Proposed Approach}

\textbf{Similarities.} The work follows the model of Anjali et al.~\cite{anjali2002} and Hamouda~\cite{hamouda2020}: an MDP with a discounted cost, a threshold policy, and the optimal policy computed by dynamic programming as a benchmark. The Markov-modulated congestion is close to Levering et al.~\cite{levering2022}, and the threshold rules to be tested are those of FlowBender~\cite{kabbani2014} and Kikuchi et al.~\cite{kikuchi2025}.

\textbf{Differences.} In the works reviewed, the switching cost is either constant, as in Hamouda~\cite{hamouda2020} and Ni\~no-Mora~\cite{ninomora2008}, or counted in changed links in a rule that is not compared with an optimum, as in Kikuchi et al.~\cite{kikuchi2025}. The works that compare with an optimum do not have a cost based on route structure, and the works that have it do not compare with an optimum. The proposed study combines the two. It also asks a different question: not whether a threshold policy is good in the scenarios of a paper, but in which conditions it fails, and it maps the failure region over the penalty, persistence, overlap and heterogeneity of the links.

\subsection{Expected Results Against the Literature}

The metrics of the comparison, and the results expected on each, are the following:
\begin{itemize}[leftmargin=*]
  \item \textbf{Optimality gap of a threshold policy.} Expected similar to the literature when the penalty is constant. Hamouda reports 1.40\% on average and 4.29\% at most~\cite{hamouda2020}, and Ni\~no-Mora below about 0.3\% for an index policy~\cite{ninomora2008}. The case study already gives a similar order: at most 0.20\% for a tuned threshold on homogeneous links with a constant penalty. Expected better than the literature where thresholds were not tuned for each scenario, and worse where the penalty depends on the route structure, since this case is not covered by the studied works.
  \item \textbf{Failure region.} The studied works do not give one, so there is no value to match. The expected contribution is the map of the conditions in which the gap exceeds a tolerance, for several tolerances.
  \item \textbf{Switching and oscillation rate.} FlowBender discusses oscillation qualitatively~\cite{kabbani2014}. The expectation is that greedy rules oscillate much more than thresholds, and that hysteresis and overlap-aware rules reduce it further, which the case study already shows for greedy (59.7\% of the switches return to the previous path).
  \item \textbf{Congestion and overhead on a real topology.} Kikuchi et al.\ measure congestion and the number of links with route changes on GEANT~\cite{kikuchi2025}. The expectation is similar congestion with a similar or lower overhead, if the overlap-aware rule is used in their setting. This is the least certain expectation and the one to be tested last.
  \item \textbf{Computation time.} The optimal policy is expected to be much slower than a threshold, as in Levering et al.~\cite{levering2022}, where the fast algorithms are introduced to avoid dynamic programming. The state space of the case study grows as $k \cdot 2^m$, so the optimal policy is only a benchmark.
\end{itemize}

\subsection{Preparing the Validation on Real Data}

The validation on a real dataset follows the setting of Kikuchi et al.~\cite{kikuchi2025}, so that the results can be compared with a published evaluation:
\begin{itemize}[leftmargin=*]
  \item \textbf{Topology and traffic:} the GEANT backbone network, with the public intradomain traffic matrices of Uhlig et al.~\cite{uhlig2006}, one matrix per 15 minutes. The $k$ shortest paths between a source and a destination are the candidate paths, with $k$ between 2 and 5.
  \item \textbf{Link states:} the traffic matrices are routed on the topology to obtain the load of every link, and a link is congested when its utilization exceeds a chosen limit. The probabilities $p$ and $q$ of each link are estimated from the sequence of utilizations, and replace the synthetic chains.
  \item \textbf{Policies:} the same as in Chapter~3, with the overhead of~\cite{kikuchi2025} (number of changed links) as one of the penalty models, next to the constant penalty and the Jaccard penalty.
  \item \textbf{Metrics:} the optimality gap, the switching and oscillation rates, and the congestion and overhead used in~\cite{kikuchi2025}.
  \item \textbf{To be confirmed:} which traffic data of~\cite{kikuchi2025} is public and in what form, and whether the links of the real topology can be modeled as independent. If the data is not available, the topology with synthetic traffic generated from the estimated parameters will be used, and this will be stated in the report.
\end{itemize}

\bibliographystyle{ieeetr}
\nocite{mckeown2008,ramalingam1996}
\bibliography{references}

\end{document}